\begin{table}[h]
\centering
\caption{Adjusting the value of the presence weight $\lambda_1$ has a minimal impact on the mean AUC over the species,  indicating that the species weights $w_s$ effectively address class imbalance issues. $\lambda_2$ is fixed at $0.8$ and $w_s$ is used. The best mean AUC in each column is highlighted in bold.}

\setlength{\tabcolsep}{0.46\tabcolsep}  % for the horizontal padding
\renewcommand{\arraystretch}{1.2} % for the vertical padding
\begin{tabular}{l|cccccc|c}
                            & AWT   & CAN   & NSW   & NZ    & SA    & SWI   & avg   \\ \hline \hline
$\lambda_1 $=0.1            & \textbf{0.709} & \textbf{0.701} & 0.701 & 0.734 & 0.784 & 0.837 & 0.744 \\
$\lambda_1 $=0.25           & 0.705 & 0.696 & 0.714 & 0.739 & 0.805 & 0.836 & 0.749 \\
$\lambda_1 $=0.5            & 0.702 & 0.695 & 0.718 & 0.740 & 0.812 & 0.836 & 0.751 \\
$\lambda_1 $=1              & 0.704 & 0.696 & \textbf{0.719} & 0.741 & \textbf{0.815} & 0.836 & \textbf{0.752} \\
$\lambda_1 $=2              & 0.704 & 0.698 & \textbf{0.719} & \textbf{0.743} & 0.814 & 0.837 & \textbf{0.752} \\
$\lambda_1 $=4              & 0.703 & 0.699 & 0.718 & \textbf{0.743} & 0.812 & 0.838 & \textbf{0.752} \\
$\lambda_1 $=10             & 0.695 & 0.696 & 0.712 & 0.737 & 0.806 & \textbf{0.840} & 0.748 \\
 \hline
\end{tabular}

\label{tab:lambda1}
\end{table}